% =============================================================================
% 2. ALGORITMOS  (30 %)
% Espacio: ~1,3 paginas (la seccion 3, tareas, vale mas: 40 %).
% El enunciado pide, por cada algoritmo, justificar POR QUE resuelve el problema.
% Cada parrafo responde dos cosas: que hace y por que funciona. El coste y donde se usa
% cada algoritmo ya estan en la Tabla de resumen: no repetirlos.
% Material para redactar: docs/algoritmos.md y docs/dificultades_encontradas.md.
% =============================================================================
\section{Algoritmos}
\label{sec:algoritmos}

\pendiente{idea de fondo (Tema 11): con una base casi ortogonal SVP y CVP son fáciles; con una
  base sesgada son difíciles. GGH publica una base mala y guarda la buena; atacarlo es conseguir
  una base buena, y eso hace LLL (Algoritmo~\ref{alg:lll} y Figura~\ref{fig:celosia}).
  La Tabla~\ref{tab:algoritmos} resume todos los algoritmos.}

% Resumen: que resuelve cada algoritmo, si da el optimo y cuanto cuesta.
% n = dimension de la celosia (o numero de incognitas), c = cota de los coeficientes.
\begin{table}[H]
\centering\small
\begin{tabularx}{\linewidth}{@{}l Y c c l@{}}
  \toprule
  Algoritmo & Qué resuelve & ¿Óptimo? & Coste & Se usa en \\
  \midrule
  Gram-Schmidt + LLL       & Base reducida (casi ortogonal) & -- & polinómico & GGH, LWE \\
  Babai (redondeo / plano) & CVP                            & no & polinómico & GGH \\
  Enumeración              & SVP y CVP                      & sí & exponencial & GGH \\
  Fuerza bruta             & SVP, CVP y LWE                 & sí & $(2c+1)^n$ & GGH, LWE \\
  Ataque primal            & Secreto de LWE                 & sí$^\ast$ & un LLL & LWE \\
  Descifrado de un bit     & Bit cifrado                    & sí$^\ast$ & lineal & LWE \\
  Sistema lineal (propio)  & Secreto de Ring-LWE            & sí$^\ast$ & $O(n^3)$ & Ring-LWE \\
  \bottomrule
\end{tabularx}
\caption{Resumen de los algoritmos. $^\ast$Si el error es lo bastante pequeño (se comprueba
  sobre el resultado).}
\label{tab:algoritmos}
\end{table}

% ---------------------------------------------------------------- Reduccion de base
\subsection{Reducción de base: Gram-Schmidt y LLL}

% LLL en pseudocodigo (igual que src/pqc/celosias/lll.py) junto a la figura de la celosia.
\begin{figure}[H]
\begin{minipage}[c]{0.58\linewidth}
  \captionof{algorithm}{LLL, como en \cite[Alg.~25]{galbraith}.}
  \label{alg:lll}
  \vspace{-4pt}\hrule\vspace{2pt}
  \begin{algorithmic}[1]
    \small
    \Require base $b_0, \dots, b_{n-1}$; $\delta = 0{,}99$
    \State $(\mu, B) \gets$ Gram-Schmidt$(b)$ \Comment{una sola vez}
    \State $k \gets 1$
    \While{$k < n$}
      \For{$j = k-1$ \textbf{hasta} $0$} \Comment{reducción de tamaño}
        \State $r \gets \lfloor \mu_{kj} \rceil$;\ \ $b_k \gets b_k - r\,b_j$
        \State actualizar la fila $k$ de $\mu$
      \EndFor
      \If{$B_k \ge (\delta - \mu_{k,k-1}^2)\,B_{k-1}$} \Comment{Lovász}
        \State $k \gets k + 1$
      \Else
        \State intercambiar $b_k$ y $b_{k-1}$; actualizar $\mu$ y $B$
        \State $k \gets \max(k - 1,\, 1)$
      \EndIf
    \EndWhile
    \State \Return $b$
  \end{algorithmic}
  \vspace{2pt}\hrule
\end{minipage}\hfill
\begin{minipage}[c]{0.38\linewidth}
  \centering
  % Base mala frente a base buena, con la celosia del Ejemplo 1 del enunciado.
%   (a) Base publica (mala):   (5, 19), (3, 14)
%   (b) Base reducida (buena): (3, 1), (-1, 4)   (LLL la devuelve con el signo cambiado)
% Las dos generan la MISMA celosia: los puntos son iguales en los dos dibujos.
% Lo unico que cambia es el paralelogramo fundamental (sombreado): largo y estrecho
% con la base mala, casi un cuadrado con la buena.
\begin{tikzpicture}[x=0.155cm,y=0.155cm,font=\footnotesize,
                    vec/.style={-{Stealth[length=3.5pt]},line width=0.8pt}]
  \newcommand{\puntos}{%
    \foreach \a in {-8,...,10}{\foreach \b in {-3,...,6}{
      \fill ({3*\a-\b},{\a+4*\b}) circle (0.38);}}}

  % (a) Base publica
  \begin{scope}
    \clip (-3,-2) rectangle (12,22);
    \fill[gray!25] (0,0) -- (5,19) -- (8,33) -- (3,14) -- cycle;
    \puntos
    \draw[vec] (0,0) -- (5,19);
    \draw[vec] (0,0) -- (3,14);
  \end{scope}
  \draw[line width=0.4pt] (-3,-2) rectangle (12,22);
  \node[anchor=west] at (5.3,19.5) {$b_1$};
  \node[anchor=east] at (2.7,14.5) {$b_2$};
  \node[below] at (4.5,-2) {(a) Base pública};

  % (b) Base reducida
  \begin{scope}[shift={(19,0)}]
    \begin{scope}
      \clip (-3,-2) rectangle (12,22);
      \fill[gray!25] (0,0) -- (3,1) -- (2,5) -- (-1,4) -- cycle;
      \puntos
      \draw[vec] (0,0) -- (3,1);
      \draw[vec] (0,0) -- (-1,4);
    \end{scope}
    \draw[line width=0.4pt] (-3,-2) rectangle (12,22);
    \node[anchor=west] at (3.2,1) {$b_1'$};
    \node[anchor=east] at (-1.1,4.3) {$b_2'$};
    \node[below] at (4.5,-2) {(b) Base reducida};
  \end{scope}
\end{tikzpicture}

  \caption{La misma celosía (Ejemplo~1 del enunciado) con la base pública y con la reducida
    por LLL.}
  \label{fig:celosia}
\end{minipage}
\end{figure}

\falta{qué hace LLL (reducción de tamaño $|\mu_{kj}| \le \frac12$ y condición de Lovász) y por qué
  funciona (solo hace restas enteras e intercambios, luego la celosía es la misma; garantía de
  longitud del primer vector). Nuestra optimización: Gram-Schmidt una sola vez y actualizar $\mu$
  y $B$ \cite[Lema~17.4.3]{galbraith}, de 329,7\,s a 0,69\,s en dimensión 31.}

% ---------------------------------------------------------------- SVP y CVP
\subsection{SVP y CVP: Babai, enumeración y fuerza bruta}

\paragraph{Babai.} \falta{redondeo: coordenadas del objetivo en la base y redondear; plano más
  cercano: capa a capa con Gram-Schmidt \cite[§18.1--18.2]{galbraith}. Por qué: acierta si
  $e\,B^{-1}$ tiene todas sus componentes en $(-\frac12, \frac12)$, por eso necesita la base
  reducida. Aproximado: en Tier02T03 da distancia$^2$ 186 y existe un punto a 155.}

\paragraph{Enumeración.} \falta{búsqueda exhaustiva acotada (Schnorr--Euchner): por qué certifica
  el óptimo y por qué es rápida tras LLL.}

\paragraph{Fuerza bruta.} \falta{probar todos los coeficientes en $[-c, c]^n$; justificar por qué
  es viable en estas tallas (lo exige el enunciado).}

% ---------------------------------------------------------------- LWE
\subsection{LWE: ataque primal y descifrado de un bit}

\paragraph{Ataque primal.} \falta{$b - As = e$ es un vector corto de una celosía construida con $A$,
  $b$ y $q$; por qué LLL lo encuentra.}

\paragraph{Descifrado de un bit.} \falta{con $s$: $R = C\cdot s$; bit 0 si $R$ está cerca de $B$,
  bit 1 si está cerca de $B + q/2$; por qué basta con que el error sea menor que $q/4$.}

% ---------------------------------------------------------------- Ring-LWE
\subsection{Ring-LWE: sistema lineal sin reducción (aportación propia)}

\falta{hipótesis: coeficientes de $a$ entre 0 y 15 y $s$, $e$ pequeños, luego $a\,s + e$ no llega a
  $q$. Idea: resolver el sistema en $\mathbb{Z}$ (matriz negacíclica de $a$) y redondear; descifrar
  con $d = c_2 - s\,c_1$. Por qué funciona: comprobación a posteriori, $e = b - a\,s$ sale en
  $\{-1, 0, 1\}$.}
